For integers $c,z$, define $c\mid z$ to mean that $z=ck$ for some
$k\in\Z$. We use this existential definition directly, including when $c=0$.
The Lean declarations below belong to \texttt{ElementaryNumberTheory.Chapter02}.

\noindent\begin{minipage}{\linewidth}
  \begin{theorem} \label{nt:thm:positive-common-divisor}
    \begin{lean}{
Suppose $a,b\in\Z$ are not both zero. There is a positive natural number $d$ which divides both $a$ and $b$, is an integer linear combination of them, and is divisible by every common divisor.
}
      theorem exists_positive_common_divisor (a b : ℤ) (hab : a ≠ 0 ∨ b ≠ 0) :
          ∃ d : ℕ, 0 < d ∧ (d : ℤ) ∣ a ∧ (d : ℤ) ∣ b ∧
            (∃ x y : ℤ, (d : ℤ) = a * x + b * y) ∧
            (∀ c : ℤ, c ∣ a → c ∣ b → c ∣ (d : ℤ))
    \end{lean}
  \end{theorem}
\end{minipage}

\begin{proof}
  \begin{lean}{
Since $a,b$ are not both zero, $a^{2}+b^{2}>0$. Thus the set of positive natural numbers of the form $ax+by$, with $x,y\in\Z$, is nonempty.
}
    classical

    have hpos : 0 < a * a + b * b := by
      rcases hab with ha | hb
      · nlinarith [sq_pos_of_ne_zero ha, sq_nonneg b]
      · nlinarith [sq_nonneg a, sq_pos_of_ne_zero hb]

    have hS : ∃ n : ℕ, 0 < n ∧ ∃ x y : ℤ, a * x + b * y = (n : ℤ) := by
      refine ⟨(a * a + b * b).toNat, ?_, a, b, ?_⟩
      · omega
      · exact (Int.toNat_of_nonneg (le_of_lt hpos)).symm
  \end{lean}

  \begin{lean}{
Let $d$ be its least element and choose $x,y\in\Z$ with $ax+by=d$. Every positive natural number represented by such a combination is at least $d$.
}
    let d : ℕ := Nat.find hS

    obtain ⟨hdpos, x, y, hxy⟩ := Nat.find_spec hS

    have hminimal (n : ℕ) (hn : 0 < n)
        (hlinear : ∃ u v : ℤ, a * u + b * v = (n : ℤ)) : d ≤ n :=
      Nat.find_min' hS ⟨hn, hlinear⟩
  \end{lean}

  \begin{lean}{
Suppose $u,v\in\Z$. Apply the division algorithm to $au+bv$ and the positive divisor $d$, obtaining $au+bv=qd+r$ with $0\le r<d$.
}
    have hddiv (u v : ℤ) : (d : ℤ) ∣ a * u + b * v := by
      obtain ⟨⟨q, r⟩, ⟨heq, hr⟩, _⟩ :=
        existsUnique_quotient_remainder (a * u + b * v) ⟨d, hdpos⟩

      change a * u + b * v = q * (d : ℤ) + (r : ℤ) at heq
      change r < d at hr
  \end{lean}

  \begin{lean}{
The remainder is again a linear combination: $r=a(u-qx)+b(v-qy)$.
}
    have hrlinear : a * (u - q * x) + b * (v - q * y) = (r : ℤ) := by
      calc
        a * (u - q * x) + b * (v - q * y) =
            (a * u + b * v) - q * (a * x + b * y) := by ring
        _ = r := by
          rw [hxy]
          linarith [heq]
  \end{lean}

  \begin{lean}{
If $r\ne0$, then $r>0$ and minimality gives $d\le r$, contradicting $r<d$. Hence $r=0$.
}
    have hrzero : r = 0 := by
      by_contra hne

      have hle := hminimal r (Nat.pos_of_ne_zero hne)
        ⟨u - q * x, v - q * y, hrlinear⟩

      omega
  \end{lean}

  \begin{lean}{
Therefore $d\mid au+bv$. Taking $(u,v)=(1,0)$ and $(0,1)$ gives $d\mid a$ and $d\mid b$.
}
      refine ⟨q, ?_⟩
      rw [hrzero] at heq
      simpa only [Nat.cast_zero, add_zero, mul_comm] using heq

    have hda : (d : ℤ) ∣ a := by
      simpa only [mul_one, mul_zero, add_zero] using hddiv 1 0

    have hdb : (d : ℤ) ∣ b := by
      simpa only [mul_one, mul_zero, zero_add] using hddiv 0 1
  \end{lean}

  \begin{lean}{
Suppose $c\mid a$ and $c\mid b$. Choose $u,v\in\Z$ with $a=cu$ and $b=cv$. Then $d=ax+by=c(ux+vy)$, so $c\mid d$. This proves all the asserted properties.
}
    have hcommon (c : ℤ) (hca : c ∣ a) (hcb : c ∣ b) : c ∣ (d : ℤ) := by
      obtain ⟨u, hu⟩ := hca
      obtain ⟨v, hv⟩ := hcb

      refine ⟨u * x + v * y, ?_⟩
      rw [← hxy, hu, hv]
      ring

    exact ⟨d, hdpos, hda, hdb, ⟨x, y, hxy.symm⟩, hcommon⟩
  \end{lean}
\end{proof}

\begin{definition}
  \begin{lean}{
For $a,b\in\Z$ not both zero, define $\gcd(a,b)$ to be the positive common divisor constructed above from the least positive linear combination. Define $\gcd(0,0)=0$.
}
    noncomputable def gcd (a b : ℤ) : ℕ :=
      if hab : a ≠ 0 ∨ b ≠ 0 then
        Classical.choose (exists_positive_common_divisor a b hab)
      else 0
  \end{lean}
\end{definition}

\noindent\begin{minipage}{\linewidth}
  \begin{theorem} \label{nt:thm:gcd-spec}
    \begin{lean}{
For integers $a,b$ not both zero, $\gcd(a,b)$ is positive, divides both inputs, is an integer linear combination of them, and is divisible by every common divisor.
}
      theorem gcd_spec (a b : ℤ) (hab : a ≠ 0 ∨ b ≠ 0) :
          0 < gcd a b ∧ (gcd a b : ℤ) ∣ a ∧ (gcd a b : ℤ) ∣ b ∧
            (∃ x y : ℤ, (gcd a b : ℤ) = a * x + b * y) ∧
            (∀ c : ℤ, c ∣ a → c ∣ b → c ∣ (gcd a b : ℤ))
    \end{lean}
  \end{theorem}
\end{minipage}

\begin{proof}
  \begin{lean}{
The definition selects the number constructed above, together with its established properties.
}
    rw [gcd, dif_pos hab]
    exact Classical.choose_spec (exists_positive_common_divisor a b hab)
  \end{lean}
\end{proof}

\noindent\begin{minipage}{\linewidth}
  \begin{theorem} \label{nt:thm:common-divisor-le-gcd}
    \begin{lean}{
Suppose $a,b\in\Z$ are not both zero and $c\in\Z$ divides both. Then $c\le\gcd(a,b)$.
}
      theorem common_divisor_le_gcd (a b c : ℤ) (hab : a ≠ 0 ∨ b ≠ 0)
          (hca : c ∣ a) (hcb : c ∣ b) : c ≤ (gcd a b : ℤ)
    \end{lean}
  \end{theorem}
\end{minipage}

\begin{proof}
  \begin{lean}{
Write $d=\gcd(a,b)>0$. Since $c\mid d$, there is $k\in\Z$ with $d=ck$. If $c>0$, positivity of $d$ forces $k>0$, hence $k\ge1$ and $d\ge c$. If $c\le0$, then $c<d$ immediately.
}
    obtain ⟨hpos, _, _, _, hcommon⟩ := gcd_spec a b hab
    have hdpos : (0 : ℤ) < gcd a b := by exact_mod_cast hpos
    obtain ⟨k, hk⟩ := hcommon c hca hcb
    by_cases hc : 0 < c
    · have hkpos : 0 < k := by nlinarith
      have hkone : 1 ≤ k := hkpos
      nlinarith
    · linarith
  \end{lean}
\end{proof}

\noindent\begin{minipage}{\linewidth}
  \begin{theorem} \label{nt:thm:gcd-zero-zero}
    \begin{lean}{
The zero-pair convention is $\gcd(0,0)=0$.
}
      theorem gcd_zero_zero : gcd 0 0 = 0
    \end{lean}
  \end{theorem}
\end{minipage}

\begin{proof}
  \begin{lean}{
Neither input is nonzero, so the definition assigns the value zero. This is a convention; the common divisors of zero and zero have no greatest integer.
}
    unfold gcd
    rw [dif_neg (by rintro (h | h); exact h rfl; exact h rfl)]
  \end{lean}
\end{proof}

\noindent\begin{minipage}{\linewidth}
  \begin{theorem} \label{nt:thm:gcd-eq-of-greatest}
    \begin{lean}{
Suppose $a,b\in\Z$ are not both zero and $d\in\N$ divides both. If every integer common divisor is at most $d$, then $\gcd(a,b)=d$.
}
      theorem gcd_eq_of_greatest (a b : ℤ) (hab : a ≠ 0 ∨ b ≠ 0) (d : ℕ)
          (hda : (d : ℤ) ∣ a) (hdb : (d : ℤ) ∣ b)
          (hgreatest : ∀ c : ℤ, c ∣ a → c ∣ b → c ≤ (d : ℤ)) : gcd a b = d
    \end{lean}
  \end{theorem}
\end{minipage}

\begin{proof}
  \begin{lean}{
Since $\gcd(a,b)$ is a common divisor, the hypothesis gives $\gcd(a,b)\le d$. The preceding bound applied to $d$ gives $d\le\gcd(a,b)$. Antisymmetry proves equality.
}
    have hspec := gcd_spec a b hab
    have hle := hgreatest (gcd a b) hspec.2.1 hspec.2.2.1
    have hge := common_divisor_le_gcd a b d hab hda hdb
    exact_mod_cast (le_antisymm hle hge)
  \end{lean}
\end{proof}

\noindent\begin{minipage}{\linewidth}
  \begin{theorem} \label{nt:thm:bezout}
    \begin{lean}{
Suppose $a,b\in\Z$ are not both zero. There are integers $x,y$ such that $\gcd(a,b)=ax+by$.
}
      theorem exists_gcd_eq_mul_add_mul (a b : ℤ) (hab : a ≠ 0 ∨ b ≠ 0) :
          ∃ x y : ℤ, (gcd a b : ℤ) = a * x + b * y
    \end{lean}
  \end{theorem}
\end{minipage}

\begin{proof}
  \begin{lean}{
The linear-combination representation is one of the properties established when constructing the gcd.
}
    exact (gcd_spec a b hab).2.2.2.1
  \end{lean}
\end{proof}

\noindent\begin{minipage}{\linewidth}
  \begin{theorem} \label{nt:thm:gcd-linear-combinations}
    \begin{lean}{
Suppose $a,b\in\Z$ are not both zero. The set of integer linear combinations $ax+by$ is exactly the set of integer multiples of $\gcd(a,b)$.
}
      theorem linear_combinations_eq_gcd_multiples (a b : ℤ) (hab : a ≠ 0 ∨ b ≠ 0) :
          {z : ℤ | ∃ x y : ℤ, z = a * x + b * y} =
            {z : ℤ | ∃ n : ℤ, z = n * (gcd a b : ℤ)}
    \end{lean}
  \end{theorem}
\end{minipage}

\enlargethispage{\baselineskip}
\begin{proof}
  \begin{lean}{
Write $d=\gcd(a,b)=ax_{0}+by_{0}$. Since $d$ divides both inputs, choose $u,v\in\Z$ with $a=du$ and $b=dv$. We prove both set inclusions.
}
    obtain ⟨x₀, y₀, hbezout⟩ := exists_gcd_eq_mul_add_mul a b hab

    have hcommon := (gcd_spec a b hab).2

    obtain ⟨u, hu⟩ := hcommon.1
    obtain ⟨v, hv⟩ := hcommon.2.1

    apply Set.ext
    intro z
    constructor
  \end{lean}

  \begin{lean}{
If $z=ax+by$, then $z=(ux+vy)d$, so $z$ is an integer multiple of $d$.
}
    · rintro ⟨x, y, hz⟩

      refine ⟨u * x + v * y, ?_⟩
      calc
        z = a * x + b * y := hz
        _ = ((gcd a b : ℤ) * u) * x + ((gcd a b : ℤ) * v) * y := by
          rw [← hu, ← hv]
        _ = (u * x + v * y) * (gcd a b : ℤ) := by ring
  \end{lean}

  \begin{lean}{
Conversely, if $z=nd$, then $z=a(nx_{0})+b(ny_{0})$, an integer linear combination of $a,b$.
}
    · rintro ⟨n, hz⟩

      refine ⟨n * x₀, n * y₀, ?_⟩
      rw [hz, hbezout]
      ring
  \end{lean}
\end{proof}
